\documentclass[11pt]{article}
\usepackage{amsmath,amssymb,amsthm}
\newtheorem{proposition}{Proposition}
\begin{document}
\title{An Exact Audit of One Two-List Tangent Calculation}
\author{}
\date{November 23, 2026}
\maketitle

\paragraph{Scope.}
This is a finite linear-algebra certificate for one specified pair of
Boolean permutation lists. It is not a universal head upper bound,
a lower bound for the full attention model, or a Lean declaration.
The earlier modular survey remains a diagnostic: equality of two ranks
modulo one prime alone would not justify equality of their real ranks.

\paragraph{The lists.}
Use coordinates $x_0,\ldots,x_{11}\in\{0,1\}$ and define
\[
 u_0(x)=\sum_{i=0}^{11}x_i,\qquad
 u_k(x)=k(1-x_{k-1})+\sum_{i=k}^{11}x_i
                      \quad(1\le k\le12).
\]
At every vertex these are a permutation of $0,\ldots,12$.
Let $v(x)=u(y)$, where
\[
 \pi=(6,0,11,8,9,1,4,10,5,2,3,7),\qquad
 y_i=x_{\pi_i}\mathbin{\mathrm{xor}}\operatorname{bit}_i(1540).
\]
Each $y_i$ is either an input coordinate or its complement, so the
$v_k$ are also affine.

Let $a_0,\ldots,a_{12},b_0,\ldots,b_{12}$ be arbitrary real affine
functions of $x$. There are $2\cdot13^2=338$ real coefficients.
Define the first variation
\[
 T_x(s)=\left.\frac{d}{d\epsilon}\right|_{\epsilon=0}
 \left[\prod_{h=0}^{12}(s+u_h(x)+\epsilon a_h(x))
       -\prod_{h=0}^{12}(s+v_h(x)+\epsilon b_h(x))\right].
\]

\begin{proposition}
If $T_x(s)$ is independent of $s$ at every Boolean vertex, then
$T_x(s)=0$ at every Boolean vertex.
\end{proposition}

\begin{proof}[Certificate proof]
Write $h(x,j)$ and $\ell(x,j)$ for the unique rows with
$u_h(x)=v_\ell(x)=j$. Put
\[
 K_{x,j}(a,b)=a_{h(x,j)}(x)-b_{\ell(x,j)}(x),\qquad
 d_j=(-1)^j j!(12-j)!.
\]
The matrix $K$ has $53248$ rows and $338$ columns, with integer
entries in $\{0,1,-1\}$. Evaluation at the distinct roots gives
\[
                         T_x(-j)=d_jK_{x,j}(a,b).
\]
Since $T_x$ has degree at most $12$, it is constant exactly when
all its $13$ displayed values agree. Define the integer matrix $C$
with $49152$ rows by
\[
 C_{x,j}=d_jK_{x,j}-d_0K_{x,0},\qquad 1\le j\le12.
\]
The hypothesis is therefore $C(a,b)=0$.

The file
\texttt{diagnostics/two-permutation-tangent-exact-n12.json}
specifies an integer matrix $V$ with $338$ rows and $15$ columns,
and a square minor of $C$ with $323$ rows. Its verifier checks:
\begin{enumerate}
\item $KV=0$ on all $4096$ Boolean vertices and all $13$ roots;
\item $V$ has column rank $15$ modulo $2147483647$;
\item the specified $323$ by $323$ integer minor of $C$ has
nonzero determinant modulo $2147483647$.
\end{enumerate}
The first check uses exact integer arithmetic. The explicit size bound
in the verifier ensures that the integer dot products cannot overflow.
The last two checks imply real independence: a nonzero determinant
modulo an integer is the residue of a nonzero integer determinant.
Consequently
\[
 \operatorname{rank}_{\mathbb R}C\ge323,\qquad
 \dim_{\mathbb R}\operatorname{span}(V)=15,
 \qquad \operatorname{span}(V)\subseteq\ker C.
\]
Rank-nullity now forces $\ker C=\operatorname{span}(V)$. Hence
$C(a,b)=0$ implies $K(a,b)=0$, so all $T_x(-j)$ vanish and
$T_x=0$.
\end{proof}

\paragraph{Reproduction.}
Run \texttt{certify\_two\_permutation\_tangent.py} with
\texttt{--verify} and the certificate path. The implementation
requires \texttt{python-flint} and \texttt{numpy}; it must be run
through the repository's \texttt{heavy} wrapper. The verifier
reconstructs the two lists and all matching equations from the
specified permutation and mask. It recomputes the minor determinant
instead of trusting the determinant value stored in the file.

The zero-output conclusion concerns first-order perturbations at
this pair. It makes no claim about another pair, a higher-order
leading term, or a general rational score away from this base.
\end{document}
